\documentclass[quiet]{ctexart}

\ExplSyntaxOn \makeatletter
\setCJKmainfont{Zhuque~Fangsong~(technical~preview)}[AutoFakeBold, ItalicFont = LXGW~WenKai]
\setCJKsansfont{LXGW~Marker~Gothic}[AutoFakeBold, AutoFakeSlant]
\setCJKmonofont{LXGW~WenKai}[AutoFakeSlant]
\ctexset{section/format+ = \raggedright}
\usepackage
  {
    geometry,mathtools,keytheorems,libertine,xeCJKfntef,
    siunitx,fixdif,derivative,physics2,tasks
  }
\usepackage[inline]{enumitem}
\usepackage{fontawesome5, fancyqr, tikz}
\hook_gput_code:nnn { cmd/faIcon/before } { . }
  {
    \group_begin:
    \fontsize { \fp_eval:n { .8 * \f@size } } { 0 } \selectfont
  }
\hook_gput_code:nnn { cmd/faIcon/after } { . } { \group_end: }
\geometry{margin = 1in}
\usepackage[warnings-off={mathtools-colon,mathtools-overbracket}]{unicode-math}
\setmathfont{STIX~Two~Math}
\usepackage[svgnames]{xcolor}
\newkeytheorem{note}
  [ style = remark, numbered = false,
    name = \color{DarkRed}\scshape\sffamily \faIcon{info-circle}~Note ]
\newkeytheorem{problem}
  [ style = definition,
    name = \color{DarkGreen}\scshape\sffamily \faIcon{pen-square}~Problem ]
\newkeytheorem{solution}
  [ style = definition, numbered = false,
    name = \color{MidnightBlue}\scshape\sffamily
           \faIcon{check-square}~Solution ]
\usephysicsmodule{ab}
\settasks { label = \sffamily\Alph*., item-indent = 2.5em, label-offset = .5pt }
\newcommand \true {\addtocounter{task}{1}\faIcon{check-circle}}
\def \@result#1 {\underline{~#1~}}
\def \s@result {\textsf{【\qquad 】}}
\DeclareRobustCommand \result {\@ifstar\s@result\@result}
\NewDocumentEnvironment{textfig}{O{7 \c_colon_str 2}}
  {
    \group_begin:
    \seq_set_split:Nee \l__whats_textfig_ratio_seq {\c_colon_str} {#1}
    \vspace{4.5pt plus 1.5pt minus 1.5pt} \noindent
    \minipage
      {
        \fp_eval:n
          {
            .96*\seq_item:Nn \l__whats_textfig_ratio_seq {1}/
              (
                \seq_item:Nn \l__whats_textfig_ratio_seq {1} + 
                \seq_item:Nn \l__whats_textfig_ratio_seq {2}
              )
          }\dimexpr \linewidth
      }
  }
  {
    \endminipage
    \vspace{4.5pt plus 1.5pt minus 1.5pt} \par
    \group_end:
  }
\NewDocumentCommand \textfigsplit {}
  {
    \endminipage
    \hspace{.02\linewidth}
    \minipage
      {
        \fp_eval:n
          {
            .96*\seq_item:Nn \l__whats_textfig_ratio_seq {2}/
              (
                \seq_item:Nn \l__whats_textfig_ratio_seq {1} + 
                \seq_item:Nn \l__whats_textfig_ratio_seq {2}
              )
          }\dimexpr \linewidth
      } 
  }

% \hook_gput_code:nnn { cmd/maketitle/before } { . } { \vskip -\baselineskip }
% \hook_gput_code:nnn { cmd/maketitle/after } { . }
%   { \__weyoung_maketitle_aftercode_aux: }
\cs_new_protected:Nn \__weyoung_maketitle_aftercode_aux:
  {
    \vspace*{-2\baselineskip}
    \begin{center}
      \begin{minipage}{.66\textwidth}
        \centering
        \begin{tasks}
          [ label-width = 5\ccwd, item-indent = 2cm,
            after-item-skip = 0pt ](2)
          \task[\textbf{课程信息}：] 理学院/A0715022
          \task[\textbf{授课教师}：] 陈许敏
          \task[\textbf{解析提供}：] \textsf{Skyrmion}
          \task[{\makebox[5em][s]{\textbf{\textsf\TeX~排版}：}}] \textsf{Skyrmion}
        \end{tasks}
      \end{minipage}
      \begin{minipage}{.15\textwidth}
        \centering
        \begin{tikzpicture}
          \node [ inner~sep = 0 ] (wechat) {
            \fancyqr
              [ size = 4em, image~padding = 1,
                image = \raisebox{4em/33}{\textcolor{DarkGreen}{\faIcon{weixin}}},
                r~color = DarkGreen, version = 4
              ]{http://weixin.qq.com/r/sRJiemTEGjllrSAZ90ej}
            };
          \node [ inner~sep = .5ex, align = center, right, font = \scriptsize ]
            at (wechat.east) {未\\央\\学\\社};
        \end{tikzpicture}
      \end{minipage}
      \begin{minipage}{.15\textwidth}
        \centering
        \begin{tikzpicture}
          \node [ inner~sep = 0 ] (qq) {
            \fancyqr
              [ size = 4em, image~x~padding = 2, image~y~padding = 1,
                image = \raisebox{4em/33}{\textcolor{MidnightBlue}{\faIcon{qq}}},
                l~color = MidnightBlue, version = 4
              ]{https://qm.qq.com/q/PH7CsQWvm2}
            };
          \node [ inner~sep = .5ex, align = center, left, font = \scriptsize ]
            at (qq.west) {物\\理\\营};
        \end{tikzpicture}
      \end{minipage}
    \end{center}
  }

\makeatother \ExplSyntaxOff
\usetikzlibrary{ fadings, patterns, arrows, decorations.markings }
\tikzset
  {
    > = stealth',
    every picture/.append style =
      { line join = round, line cap = round, thick }
  }

\title{\sffamily\bfseries 2021年HDU「高等数学A1」期中试题解答}
\author{Mingyu Xia}

\begin{document}

\maketitle

\section{选择题（每题3分，共27分）}

\begin{problem}
  若 $f(x)$ 在 $x_0$ 的导函数 $f'(x_0)$ 存在，则
  $\lim\limits_{n \to \infty} \ab[f\ab(x_0 + \frac1n) - f\ab(x_0 - \frac1n)] = $
  \begin{tasks}(4)
    \task $f'(x_0)$
    \task[\true] $2f'(x_0)$
    \task $3f'(x_0)$
    \task $4f'(x_0)$
  \end{tasks}
\end{problem}
\begin{solution}

\end{solution}

\begin{problem}
  设函数 $f(x)$ 在 $[0,1]$ 上有 $f''(x) > 0$，则有不等式
  \begin{tasks}(2)
    \task $f'(1) > f'(0) > f(1) - f(0)$
    \task $f'P(0) > f(1) - f(0) > f'(1)$
    \task $f(1) - f(0) > f'(1) > f'(0)$
    \task[\true] $f'(1)$
  \end{tasks}
\end{problem}
\begin{solution}
  
\end{solution}

\begin{problem}
  设 $f(x)$，$g(x)$ 有连续导函数，则下列命题中，正确的命题一共有
  \begin{tasks}[label = \roman*.](6)
    \task*(3) 若 $\int f(x) \d x = g(x)$，则 $f(x) = g'(x)$
    \task*(3) 若 $f'(x) = g(x)$，则 $f(x) = \int g(x) \d x$
    \task*(2) $\int (f(x))' \d x = f(x)$
    \task*(2) $(\int f(x) \d x)' = f(x)$
    \task*(2) $f(x) \d x = \d[\int f(x) \d x]$
  \end{tasks}
  \begin{tasks}(4)
    \task 3个
    \task[\true] 2个
    \task 1个
    \task 0个
  \end{tasks}
\end{problem}
\begin{solution}
  
\end{solution}

\begin{problem}
  设 $f(x) = \int_0^x \sin^2 t \d t$，
  则当 $x \to 0$，$f(x)$ 是 $x^2$ 的\result*无穷小
  \begin{tasks}(4)
    \task 等价
    \task[\true] 高阶
    \task 同阶，但非等价
    \task 低阶
  \end{tasks}
\end{problem}
\begin{solution}
  
\end{solution}

\begin{problem}
  已知 $f'(\cos^2 x) = \sin^2 x$，则 $f(x) = $
  \begin{tasks}(4)
    \task $x - \frac12 x^2 + C$
    \task[\true] $x + \frac12 x^2 + C$
    \task $x^2 + \frac12 x + C$
    \task $x^2 - \frac12 x + C$
  \end{tasks}
\end{problem}
\begin{solution}
  
\end{solution}

\begin{problem}
  关于反常积分 $\int_2^{+\infty} \frac{\d x}{(x - 1)^p}$ 敛散性下列结论正确的是
  \begin{tasks}(4)
    \task[\true] $p < 1$ 收敛
    \task $p > 1$ 收敛
    \task $p \in \mathbb R$ 都发散
    \task $p \in \mathbb R$ 都收敛
  \end{tasks}
\end{problem}
\begin{solution}
  
\end{solution}

\begin{problem}
  由 $y_1 = e^x$，$y_2 = e^{-x}$ 及 $x = 1$ 所围成的平面图区域的面积表示为
  \begin{tasks}(4)
    \task[\true] $\int_0^1 (e^x - e^{-x}) \d x$
    \task $\int_0^e (e^x - e^{-x}) \d x$
    \task $\int_{e^{-1}}^e (1 - \ln y) \d y$
    \task $2\int_{e^{-1}}^e (1 - \ln y) \d y$
  \end{tasks}
\end{problem}
\begin{solution}
  
\end{solution}

\begin{problem}
  对数螺线 $\rho = e^\theta$（或记为 $\theta  = \ln\rho$）
  相应于 $\theta$ 从 $0$ 到 $\pi$ 的弧长为
  \begin{tasks}(4)
    \task $\sqrt2(\pi - 1)$
    \task $2(e^\pi - 1)$
    \task[\true] $\sqrt2(e^\pi - 1)$
    \task $2(\pi - 1)$
  \end{tasks}
\end{problem}
\begin{solution}
  
\end{solution}

\begin{problem}
  设连续函数 $f(x)$ 满足 $\int_0^x f(t) \d t = f(x) - 1$，则 $f(x) = $
  \begin{tasks}(4)
    \task $y = 2e^x$
    \task $y = e^{2x}$
    \task $y  e^x$
    \task $y = e^{x^2}$
  \end{tasks}
\end{problem}
\begin{solution}
  
\end{solution}

\section{填空题}

\begin{problem}
  极限 $\lim\limits_{x \to 0} \ab(\frac{1}{x\tan x} - \frac{1}{x^2}) = $
  \result{$-\frac13$}.
\end{problem}
\begin{solution}
  
\end{solution}

\begin{problem}
  $\int_{-2021}^{2021} \frac{x\sqrt{1 + \sin^2 x}}{1 + x^2} \d x = $
  \result{$0$}.
\end{problem}
\begin{solution}
  
\end{solution}

\begin{problem}
  已知 $f(x) =
  \begin{cases}
    1 - \arctan \frac1x, & x < 0\\
    a + e^{2\sqrt x},    & x \geq 0
  \end{cases}$
  是连续函数，则 $a = $ \result{$\frac\pi2$}.
\end{problem}
\begin{solution}
  
\end{solution}

\section{解答题}

\begin{problem}
  计算 $\int \frac{2x - 3}{x^2 - 2x + 3} \d x$.
\end{problem}
\begin{solution}
  \[\ln(x^2 -2x + 3) - \frac{1}{\sqrt2} \arctan \frac{x - 1}{\sqrt2} + C\]
\end{solution}

\begin{problem}
  设隐函数 $y = y(x)$ 由方程 $e^y + \sin(xy) - \int_0^x e^t \d t = e$
  所确定，求 $y'(0)$.
\end{problem}
\begin{solution}
  \[0\]
\end{solution}

\begin{problem}
  已知 $\int \frac{f(x)}{\sqrt{1 + 2x - x^2}} \d x = x + C$，
  求 $\int \frac{1}{f(x)} \d x$.
\end{problem}
\begin{solution}
  \[\arcsin \frac{x - 1}{\sqrt 2} + C\]
\end{solution}

\begin{problem}
  求微分方程 $\odv yx = e^{x - y}$ 满足初始条件 $y(0) = 1$ 的特解.
\end{problem}
\begin{solution}
  \[e^y = e^x + e - 1\]
\end{solution}

\begin{problem}
  求曲线 $f(x) = \int_0^{x^2} \ab(\frac32 - \frac12t) \d t$ 的凹凸区间和拐点.
\end{problem}
\begin{solution} \leavevmode
  \begin{itemize}
    \item 凸区间 $(-\infty, -1]$ 与 $[1, +\infty)$，凹区间 $[-1,1]$
    \item 拐点 $(-1,5/4)$ 与 $(1,5/4)$.
  \end{itemize}
\end{solution}

\begin{problem}
  已知 $ \lim\limits_{x \to \infty} \ab(1 - \frac1x)^{-2ax}
      = \int_{-\infty}^a xe^{2x} \d x$，求常数 $a$.
\end{problem}
\begin{solution}
  \[\frac52\]
\end{solution}

\begin{problem}
  求定积分 $\int_1^3 f(x - 2) \d x$，其中 $f(x) =
  \begin{cases}
    xe^{x^2},             & x < 0\\
    \frac{1}{1 + \cos x}, & x \geq 0
  \end{cases}$.
\end{problem}
\begin{solution}
  \[\frac12 - \frac12e + \tan\frac12\]
\end{solution}

\begin{problem}
  求微分方程 $y'' - 5y' + 6y = xe^{2x}$ 的通解.
\end{problem}
\begin{solution}
  \[y = C_1e^{2x} + C_2e^{3x} - \ab(\frac12x^2 + x)e^{2x}\]
\end{solution}

\begin{problem}
  过坐标原点作曲线 $y = \ln x$ 的切线, 
  该切线与曲线 $y = \ln x$ 及 $x$ 轴围成的平面图形 $D$. 试求
  \begin{enumerate}
    \item 曲线 $y = \ln x$ 过坐标原点的切线方程.
    \item 平面图形 $D$ 绕 $y$ 轴旋转一周所得的旋转体的体积.
  \end{enumerate}
\end{problem}
\begin{solution}
  \begin{enumerate*}
    \item $y = \frac x3$.
    \item $\ab(\frac16e^2 - \frac12)\pi$.
  \end{enumerate*}
\end{solution}

\begin{problem}
  设 $a > 0$，求证
  $ \int_0^\pi xa^{\sin x} \d x \cdot
    \int_0^{\pi/2} a^{-\cos x} \d x \geq \frac{\pi^3}{4}$.
\end{problem}
\begin{solution}
  \begin{proof}
    
  \end{proof}
\end{solution}

\end{document}